On considère l’équation différentielle (E) : $y' + 2y = 3x+2$ où $y$ désigne une fonction de la variable réelle $x$ définie et dérivable sur $\R$, $y'$ sa fonction dérivée.

\begin{enumerate}
	\item Donner, puis résoudre l'équation homogène (E$_0$) :\sautp

	\Lignes{2}
	\item A l’aide du logiciel Géogébra, déterminer les nombres réels $a$ et $b$ pour que la fonction $p$ définie sur $\R$ par : $p(x) = ax+b$ soit une solution particulière de l'équation différentielle (E), et donner l’expression de $p$. 

	
	On trouve : \hfill $a=$\makebox[0.1\linewidth]{\dotfill} \hfill $b=$\makebox[0.1\linewidth]{\dotfill} \hfill $p(x)=$\makebox[0.3\linewidth]{\dotfill}

	\smallskip
	
	\appelprof
	
	\item En déduire l'ensemble des solutions de (E) : \sautp

	\Lignes{1}
	\item Déterminer, par un calcul, la solution $f$ de (E) telle que $f(0)=2$ : \sautp

	\Lignes{4}

\end{enumerate}